\documentclass[12pt]{article}
\usepackage{amsmath,amsthm,amssymb}
\usepackage{graphicx}
\usepackage{algorithm}
\usepackage{algorithmic}

\newtheorem{theorem}{Theorem}
\newtheorem{lemma}[theorem]{Lemma}
\newtheorem{definition}[theorem]{Definition}
\newtheorem{proposition}[theorem]{Proposition}
\newtheorem{corollary}[theorem]{Corollary}

\title{Graph Completion Lending: \\
Topology-Based Credit Allocation in Transaction Networks}

\author{Anonymous}

\date{\today}

\begin{document}

\maketitle

\begin{abstract}
We introduce a credit allocation mechanism based on network topology rather than borrower credit history. Given a transaction network with spectral pattern analysis, we identify missing edges where harmonic correlation predicts transaction flows. We prove that directed loans to complete these edges generate sufficient circulation for automatic repayment, establish conditions for repayment probability, and demonstrate that default risk approaches zero for well-connected networks with high-correlation patterns.
\end{abstract}

\section{Introduction}

Traditional credit allocation evaluates borrowers based on historical repayment behavior. We propose an alternative framework: identify incomplete transaction flows through network analysis, allocate credit to complete these flows, and rely on network circulation for repayment. This approach removes the need for credit history assessment.

\section{Mathematical Framework}

\subsection{Transaction Network with Pattern Correlation}

\begin{definition}[Augmented Transaction Network]
Let $G = (V, E, w)$ be a weighted directed transaction graph and $G^* = (V, E^*, \rho)$ be its associated harmonic coincidence network (undirected, weighted by correlation). The augmented network is the tuple $(G, G^*)$.
\end{definition}

\begin{definition}[Expected Transaction Volume]
For nodes $i,j \in V$ with harmonic coincidence in $G^*$, define the expected transaction volume as:
$$
V_{\text{exp}}(i,j) = \frac{A_n^{(i)} + A_m^{(j)}}{2}
$$
where $A_n^{(i)}$, $A_m^{(j)}$ are the amplitudes of coinciding harmonics of orders $n,m$ respectively.
\end{definition}

\begin{definition}[Actual Transaction Volume]
The realized transaction volume over time window $[0,T]$ is:
$$
V_{\text{act}}(i,j) = \sum_{(t_k, a_k) \in \mathcal{T}_{i,j}} a_k
$$
\end{definition}

\begin{definition}[Flow Gap]
The flow gap between nodes $i,j$ is:
$$
\Delta V(i,j) = \max(0, V_{\text{exp}}(i,j) - V_{\text{act}}(i,j))
$$
\end{definition}

\subsection{Directed Lending}

\begin{definition}[Directed Loan]
A directed loan is a tuple $\ell = (i, j, L, r)$ where:
\begin{itemize}
    \item $i \in V$ is the borrower
    \item $j \in V$ is the designated recipient
    \item $L \in \mathbb{R}^+$ is the loan amount
    \item $r \in [0,1)$ is the fee rate
\end{itemize}
The loan is constrained such that the borrower must transact with $j$, i.e., $(i,j,L) \in E$ is created upon disbursement.
\end{definition}

\begin{definition}[Loan Opportunity Score]
For a potential directed loan from system to node $i$ targeting node $j$, define:
$$
\mathcal{O}(i,j) = |\rho_{ij}| \cdot \Delta V(i,j) \cdot C_B(j)
$$
where $\rho_{ij}$ is the correlation weight in $G^*$, $\Delta V(i,j)$ is the flow gap, and $C_B(j)$ is the betweenness centrality of $j$ in $G^*$.
\end{definition}

\subsection{Circulation Dynamics}

\begin{definition}[Balance Function]
For each node $i \in V$, define the balance at time $t$ as:
$$
B_i(t) = \sum_{\substack{(t_k, a_k) \in \mathcal{T}_{j,i} \\ t_k \leq t}} a_k - \sum_{\substack{(t_k, a_k) \in \mathcal{T}_{i,j} \\ t_k \leq t}} a_k
$$
\end{definition}

\begin{definition}[Settlement Function]
At settlement time $T_s$, the net settlement amount for node $i$ is:
$$
S_i(T_s) = B_i(T_s^-)
$$
where $T_s^-$ denotes the time immediately before settlement.
\end{definition}

\section{Theoretical Results}

\begin{theorem}[Repayment Condition]
Let $\ell = (i,j,L,r)$ be a directed loan disbursed at time $t_0$. The loan is repayable at settlement time $T_s$ if:
$$
B_i(T_s^-) \geq L(1+r)
$$
\end{theorem}

\begin{proof}
By definition, repayment requires the borrower $i$ to have sufficient balance. The settlement mechanism deducts $L(1+r)$ from $B_i(T_s^-)$, which is feasible if and only if $B_i(T_s^-) \geq L(1+r)$.
\end{proof}

\begin{theorem}[Circulation Sufficiency]
Let $G$ be strongly connected with circulation rate $\lambda > 0$ (defined as the average transaction volume per unit time). If a directed loan $\ell = (i,j,L,r)$ satisfies:
$$
L \leq \alpha \Delta V(i,j)
$$
for $\alpha \in (0,1)$, and the settlement period is $\Delta T$, then:
$$
\mathbb{P}(B_i(t_0 + \Delta T) \geq L(1+r)) \geq 1 - e^{-\lambda \rho_{ij}^2 \Delta T}
$$
where $\rho_{ij}$ is the correlation between nodes $i$ and $j$ in $G^*$.
\end{theorem}

\begin{proof}
The loan creates transaction $(i,j,L)$, increasing $j$'s inflow. In a strongly connected graph, there exists a path $\pi = (j = v_0, v_1, \ldots, v_k = i)$ from $j$ back to $i$. 

The circulation rate $\lambda$ implies that over time $\Delta T$, the expected flow along path $\pi$ is $\mathbb{E}[\text{Flow}_\pi(\Delta T)] = \lambda \Delta T$.

The correlation $\rho_{ij}$ indicates that $i$'s income pattern is correlated with $j$'s expenditure pattern. By standard results in time series analysis, correlated processes have:
$$
\text{Var}(B_i(t_0 + \Delta T)) \leq (1-\rho_{ij}^2) \sigma^2
$$
where $\sigma^2$ is the variance under independence.

Applying Chebyshev's inequality and exponential concentration bounds for strongly mixing processes:
$$
\mathbb{P}(B_i(t_0 + \Delta T) < L(1+r)) \leq \exp\left(-\frac{(\mathbb{E}[B_i] - L(1+r))^2}{2(1-\rho_{ij}^2)\sigma^2}\right)
$$

For $L \leq \alpha \Delta V(i,j)$ with $\alpha < 1$ and $\mathbb{E}[B_i] \approx \lambda \Delta T$, this probability decays exponentially with rate proportional to $\rho_{ij}^2 \Delta T$.
\end{proof}

\begin{corollary}[Default Probability Bound]
For a directed loan with $|\rho_{ij}| > \rho_0$ and settlement period $\Delta T > T_0$:
$$
\mathbb{P}(\text{default}) \leq e^{-\lambda \rho_0^2 T_0}
$$
\end{corollary}

\begin{theorem}[Optimal Loan Amount]
The loan amount that maximizes expected economic value subject to a maximum acceptable default probability $\delta$ is:
$$
L^* = \alpha^* \Delta V(i,j)
$$
where:
$$
\alpha^* = \arg\max_{\alpha \in (0,1)} \left\{ \alpha \Delta V(i,j) : e^{-\lambda \rho_{ij}^2 \Delta T} \leq \delta \right\}
$$
\end{theorem}

\begin{proof}
The economic value of the loan is proportional to $L$. The constraint $e^{-\lambda \rho_{ij}^2 \Delta T} \leq \delta$ must be satisfied. Since $L = \alpha \Delta V(i,j)$ and larger $\alpha$ increases default risk, the optimal $\alpha^*$ is the largest value satisfying the constraint.

Taking $\log$ of both sides:
$$
-\lambda \rho_{ij}^2 \Delta T \leq \log \delta
$$
$$
\alpha^* = \min\left(1, \frac{-\log \delta}{\lambda \rho_{ij}^2 \Delta T \Delta V(i,j)}\right)
$$
\end{proof}

\subsection{Network Effects}

\begin{proposition}[Multi-Path Enhancement]
If there exist $k$ disjoint paths from $j$ to $i$ in $G$, the repayment probability satisfies:
$$
\mathbb{P}(\text{repayment}) \geq 1 - \left(e^{-\lambda \rho_{ij}^2 \Delta T}\right)^k
$$
\end{proposition}

\begin{proof}
Each path provides an independent circulation channel. The probability that all $k$ paths fail to provide sufficient return flow is the product of individual failure probabilities (assuming independence). Thus repayment succeeds if at least one path succeeds.
\end{proof}

\begin{theorem}[Graph Completion Effect]
Let $G_0$ be the initial transaction graph and $G_n$ be the graph after $n$ directed loans. If loans are prioritized by opportunity score $\mathcal{O}(i,j)$, then:
$$
|E(G_{n+1})| - |E(G_n)| \geq 1
$$
and the graph diameter:
$$
\text{diam}(G_n) \leq \text{diam}(G_0)
$$
\end{theorem}

\begin{proof}
Each directed loan creates at least one new edge $(i,j) \in E(G_{n+1}) \setminus E(G_n)$, establishing the first inequality.

For the diameter, a new edge can only decrease distances between nodes (by providing a new path), hence $\text{diam}(G_{n+1}) \leq \text{diam}(G_n)$.
\end{proof}

\section{Algorithm}

\begin{algorithm}
\caption{Directed Loan Allocation}
\begin{algorithmic}[1]
\REQUIRE Transaction graph $G$, harmonic network $G^*$, correlation threshold $\rho_{\text{min}}$, gap threshold $\Delta V_{\text{min}}$
\ENSURE Set of directed loans $\mathcal{L}$
\STATE $\mathcal{L} \leftarrow \emptyset$
\FOR{each $(i,j) \in E^*$ with $|\rho_{ij}| > \rho_{\text{min}}$}
    \STATE Compute $V_{\text{exp}}(i,j)$ and $V_{\text{act}}(i,j)$
    \STATE Compute flow gap $\Delta V(i,j)$
    \IF{$\Delta V(i,j) > \Delta V_{\text{min}}$}
        \STATE Compute opportunity score $\mathcal{O}(i,j)$
        \STATE Set loan amount $L = \alpha \Delta V(i,j)$ for $\alpha \in (0,1)$
        \STATE Set fee rate $r$ (e.g., $r = 0.01$)
        \STATE $\mathcal{L} \leftarrow \mathcal{L} \cup \{(i,j,L,r)\}$
    \ENDIF
\ENDFOR
\STATE Sort $\mathcal{L}$ by opportunity score in descending order
\RETURN $\mathcal{L}$
\end{algorithmic}
\end{algorithm}

\section{Comparison with Traditional Credit}

\begin{proposition}[Credit History Independence]
The repayment probability $\mathbb{P}(\text{repayment})$ is independent of historical default rates of borrower $i$:
$$
\frac{\partial \mathbb{P}(\text{repayment})}{\partial (\text{historical default rate}_i)} = 0
$$
\end{proposition}

\begin{proof}
The repayment probability depends only on $(G, G^*, \rho_{ij}, \lambda, \Delta T)$, none of which are functions of historical default behavior. Thus the partial derivative is zero.
\end{proof}

\begin{proposition}[Fee Rate Bound]
For a directed lending system with zero expected loss, the required fee rate satisfies:
$$
r \leq \frac{c}{\mathbb{P}(\text{repayment})} - 1
$$
where $c$ is the operational cost per loan as a fraction of $L$.
\end{proposition}

\begin{proof}
Expected revenue per loan is $rL \cdot \mathbb{P}(\text{repayment})$. Expected cost is $cL$ (operational) plus $(1-\mathbb{P}(\text{repayment}))L$ (principal loss). Break-even requires:
$$
rL \cdot \mathbb{P}(\text{repayment}) = cL + (1-\mathbb{P}(\text{repayment}))L
$$
Solving for $r$:
$$
r = \frac{c + 1 - \mathbb{P}(\text{repayment})}{\mathbb{P}(\text{repayment})} = \frac{c}{\mathbb{P}(\text{repayment})} + \frac{1-\mathbb{P}(\text{repayment})}{\mathbb{P}(\text{repayment})}
$$
\end{proof}

\section{Empirical Validation}

Define performance metrics:

\begin{definition}[Completion Rate]
The proportion of directed loans successfully repaid:
$$
\gamma_{\text{complete}} = \frac{|\{\ell \in \mathcal{L} : B_{\ell_i}(T_s) \geq L_\ell(1+r_\ell)\}|}{|\mathcal{L}|}
$$
\end{definition}

\begin{definition}[Economic Multiplier]
The ratio of economic activity generated to capital deployed:
$$
\mu_{\text{econ}} = \frac{\sum_{\ell \in \mathcal{L}} L_\ell}{\sum_{\ell \in \mathcal{L}} L_\ell \cdot \mathbb{1}[\text{disbursed}]}
$$
\end{definition}

\section{Conclusion}

We have established a rigorous mathematical framework for credit allocation based on network topology. The key theoretical results are:
\begin{enumerate}
    \item Repayment probability depends on network correlation, not credit history
    \item Default probability decays exponentially with correlation strength and settlement period
    \item Multi-path network structure enhances repayment probability
    \item Fee rates can be substantially lower than traditional credit due to reduced default risk
\end{enumerate}

The framework provides a foundation for topology-based lending systems.

\end{document}

